\documentclass{article}
\usepackage[T1]{fontenc}
\usepackage{lmodern}
\usepackage{graphicx}
\usepackage{amsmath,amssymb}
\usepackage[a4paper,margin=2cm]{geometry}
\usepackage[absolute,overlay]{textpos}
\setlength{\TPHorizModule}{1mm}
\setlength{\TPVertModule}{1mm}
\usepackage[indonesian]{babel}
\usepackage{hyperref}
\setlength{\emergencystretch}{2em}
% unit: Halaman 18

\begin{document}

% === Halaman 18 (llm) ===

Blok plainteks ke-1 : 011000 \\
Kunci publik : 62, 93, 81, 88, 102, 37 \\
Kriptogram : $(1 \times 93) + (1 \times 81) = 174$

\vspace{5mm}

Blok plainteks ke-2 : 110101 \\
Kunci publik : 62, 93, 81, 88, 102, 37 \\
Kriptogram : $(1 \times 62) + (1 \times 93) + (1 \times 88) + (1 \times 37) = 280$

\vspace{5mm}

Blok plainteks ke-3 : 101110 \\
Kunci publik : 62, 93, 81, 88, 102, 37 \\
Kriptogram : $(1 \times 62) + (1 \times 81) + (1 \times 88) + (1 \times 102) = 333$

\vspace{5mm}

Jadi, cipherteks yang dihasilkan : 174, 280, 333

\end{document}
